% Solutions/hints for ch23 exercises (assembled into Appendix C)
\begin{solution}{exr:ch23-laplacian}
With the edges $\{1,2\},\{2,3\},\{3,4\},\{4,1\},\{1,3\}$ the degrees are $(3,2,3,2)$ and
\[
\mat{L}=\begin{pmatrix} 3&-1&-1&-1\\ -1&2&-1&0\\ -1&-1&3&-1\\ -1&0&-1&2\end{pmatrix}.
\]
Every row sums to zero, so $\mat{L}\vect{1}=\vect{0}$. For $\vect{x}=(1,0,0,1)\T$ the quadratic form is $\vect{x}\T\mat{L}\vect{x}=\sum_{\{i,j\}\in E}(x_i-x_j)^2=1+0+1+0+1=3$ (the edges $\{1,2\}$, $\{3,4\}$ and $\{1,3\}$ each contribute $1$). Guess eigenvectors orthogonal to $\vect{1}$: $\mat{L}(0,1,0,-1)\T=2\,(0,1,0,-1)\T$, $\mat{L}(1,0,-1,0)\T=4\,(1,0,-1,0)\T$ and $\mat{L}(1,-1,1,-1)\T=4\,(1,-1,1,-1)\T$, so the spectrum is $\{0,2,4,4\}$ (check: the trace is $10$). Hence $\lambda_2=2$, the time constant is $1/\lambda_2=0.5$~s, the safe step size is $\eps<1/d_{\max}=1/3$ and the exact bound is $\eps<2/\lambda_4=1/2$. A useful shortcut: the complement of the graph in $K_4$ is the single edge $\{2,4\}$ with Laplacian spectrum $\{0,0,0,2\}$, and $\mat{L}(G)+\mat{L}(\bar G)=4\mat{I}-\vect{1}\vect{1}\T$ gives the non-zero eigenvalues of $G$ as $4-\{2,0,0\}$.
\end{solution}

\begin{solution}{exr:ch23-proof}
Write $\mat{P}=\mat{I}-\eps\mat{L}$. It is symmetric, and if $\mat{L}\vect{v}_k=\lambda_k\vect{v}_k$ then $\mat{P}\vect{v}_k=(1-\eps\lambda_k)\vect{v}_k$, so $\mat{P}$ has the eigenvalues $\mu_k=1-\eps\lambda_k$ with the same orthonormal eigenvectors. Because $\vect{1}\T\mat{L}=\vect{0}\T$ we have $\vect{1}\T\vect{x}_{k+1}=\vect{1}\T\vect{x}_k$: the average $\bar x$ is preserved. Let $\vect{\delta}_k=\vect{x}_k-\bar x\vect{1}$; then $\vect{\delta}_{k+1}=\mat{P}\vect{\delta}_k$ and $\vect{\delta}_k\perp\vect{1}$ for all $k$. Expanding $\vect{\delta}_0=\sum_{k\ge2}c_k\vect{v}_k$ gives $\vect{\delta}_m=\sum_{k\ge2}c_k\mu_k^m\vect{v}_k$ and, by orthonormality, $\norm{\vect{\delta}_m}\le\rho^m\norm{\vect{\delta}_0}$ with $\rho=\max_{k\ge2}\abs{1-\eps\lambda_k}=\max(\abs{1-\eps\lambda_2},\abs{1-\eps\lambda_n})$. It remains to show $\rho<1$: since the graph is connected, $\lambda_2>0$, and $\eps>0$ gives $1-\eps\lambda_k<1$ for $k\ge2$; since $\lambda_n\le 2d_{\max}$ (Gershgorin: every eigenvalue of $\mat{L}$ lies in a disc centred at some $d_i$ with radius $d_i$) and $\eps<1/d_{\max}$, we get $1-\eps\lambda_k\ge1-\eps\lambda_n>1-2=-1$. So every $\mu_k$, $k\ge2$, lies strictly inside $(-1,1)$, $\rho<1$, and $\vect{x}_m\to\bar x\vect{1}$ geometrically. The exact condition is $\eps<2/\lambda_n$; the bound $1/d_{\max}$ only uses information every drone has locally.
\end{solution}

\begin{solution}{exr:ch23-inconsistent}
(a) Summing the protocol over $i$ gives $\sum_i\dot{\pos}_i=\sum_i\sum_j a_{ij}(\pos_j-\pos_i)-\sum_i\sum_j a_{ij}\vect{d}_{ij}$. The first double sum vanishes because the graph is undirected (each pair contributes $(\pos_j-\pos_i)+(\pos_i-\pos_j)$). The second vanishes if and only if $\vect{d}_{ji}=-\vect{d}_{ij}$ on every edge; otherwise the centroid moves with the constant velocity $-\frac1n\sum_i\sum_j a_{ij}\vect{d}_{ij}$ and the swarm drifts away for ever, even though the relative positions settle. (b) With antisymmetric $\vect{d}_{ij}$ the protocol is $\dot{\pos}_i=-\nabla_{\pos_i}J$ for $J(\pos)=\tfrac12\sum_{\{i,j\}\in E}\norm{\pos_j-\pos_i-\vect{d}_{ij}}^2$ (differentiate one term with respect to $\pos_i$ to check). $J$ is a convex quadratic bounded below, so gradient descent converges to a minimiser and $J$ decreases monotonically. If the vectors close around every cycle, $J_{\min}=0$ and the minimisers are the translated formations; otherwise $J_{\min}>0$ and the swarm settles into a least-squares compromise with residual formation error $e_F=\sqrt{2J_{\min}/\abs{E}}$. For the triangle with $\vect{d}_{12}=\vect{d}_{23}=(1,0)\T$ and $\vect{d}_{31}=(-1,0)\T$ the mismatch $(1,0)\T$ is spread evenly, each edge ends with error $(1/3,0)\T$ and $e_F=1/3$~m. (c) Storing offsets $\vect{o}_i$ and computing $\vect{d}_{ij}=\vect{o}_j-\vect{o}_i$ makes both antisymmetry and cycle closure automatic.
\end{solution}

\begin{solution}{exr:ch23-secondorder}
(a) In the eigenbasis of $\mat{L}$ the mode with eigenvalue $\lambda$ obeys $\ddot q+(k_v\lambda+k_d)\dot q+k_p\lambda q=0$, with characteristic polynomial $s^2+(k_v\lambda+k_d)s+k_p\lambda$. A quadratic with positive coefficients has both roots in the open left half-plane, so every mode with $\lambda>0$ decays if $k_p>0$ and $k_v\lambda+k_d>0$; the mode $\lambda=0$ gives $s\in\{0,-k_d\}$. With $k_v=k_d=0$ the roots are $s=\pm\mathrm{i}\sqrt{k_p\lambda}$: undamped oscillation at the angular frequencies $\sqrt{k_p\lambda_k}$, which for the example graph and $k_p=1$ are $1$, $\sqrt3$ and $2$~rad/s. For $k_d=0$ the damping ratio of mode $\lambda$ is $\zeta=\tfrac12 k_v\sqrt{\lambda/k_p}$; choosing $k_v=2\sqrt{k_p/\lambda_2}$ makes the slowest mode critically damped and the faster ones overdamped. (b) Without feed-forward the error $\vect{y}=\pos-\vect{o}-\vect{r}(t)$ of the first-order leader-follower controller obeys $\dot{\vect{y}}=-(k\mat{L}+k_l\mat{B})\vect{y}-\dot{\vect{r}}\vect{1}$, so with $\dot{\vect{r}}=v\,\vect{e}$ constant the steady state is $\vect{y}_{ss}=-v\,(k\mat{L}+k_l\mat{B})\inv\vect{1}\otimes\vect{e}$. For the example graph, $k=k_l=1$ and drone~1 pinned, solving $(\mat{L}+\mat{B})\vect{z}=\vect{1}$ gives $\vect{z}=(4,\,16/3,\,17/3,\,20/3)\T$: at $v=1$~m/s the leader trails its reference by $4$~m and the tail drone by $6.67$~m, and the edge errors $z_j-z_i$ give $e_F=\sqrt{(16/9+25/9+1/9+1)/4}\approx1.19$~m for ever. Adding $\dot{\vect{r}}$ to every drone's command removes the forcing term and the lag.
\end{solution}
